\section{Focused cyclic sequent calculus}
\label{sec:focused}

A sequent has the form $\Gamma \vdash t : A$, representing the goal of finding a term $t$ of type $A$ in context $\Gamma$. A \emph{cyclic proof} is a finite directed graph where each vertex $v$ is labelled with a sequent, has successors $\succof(v)$ (its premises), and satisfies a local rule. The graph may contain cycles. A proof is valid if every vertex's local rule holds and every infinite path contains infinitely many \emph{progress events} (splits).

\subsection{Asynchronous rules (driving)}

The asynchronous rules are deterministic and invertible. They correspond to CBN reduction: $\beta$-reduction, fixpoint unfolding, $\iota$-reduction (case-of-constructor), and commuting conversions. All have the same shape---a single premise that is the result of one computation step---and are applied eagerly without backtracking.

\subsection{Synchronous rule: case split}

The synchronous rule introduces choice by analyzing a neutral variable:

\begin{center}
\AxiomC{$\Gamma \vdash b_0 : C[\Roll{I}{0}{}/]$}
\AxiomC{$\cdots$}
\AxiomC{$\Gamma \vdash b_{n-1} : C[\Roll{I}{n-1}{}/]$}
\TrinaryInfC{$\Gamma, x : \Ind{I}{\vec{a}} \vdash \Case{I}{x}{C}{\vec{b}} : C[x/]$}
\DisplayProof
\quad (\textsc{split-case-var})
\end{center}

Each premise corresponds to one constructor of the inductive type. \textbf{Key property:} this is a \emph{progress event}. Every cycle must pass through at least one split, ensuring structural descent.

\subsection{Backlink rules (folding)}

A backlink closes a goal by pointing to a previous vertex. This is the cyclic aspect: rather than choosing an induction scheme upfront, the backlink \emph{is} the induction hypothesis, discovered when a current goal is recognized as an instance of a previous one. The trace condition ensures the resulting recursion is well-founded.

\begin{center}
\AxiomC{}
\UnaryInfC{$\Gamma \vdash t : A$}
\DisplayProof
\quad (\textsc{back-exact})
\qquad
\AxiomC{$\Delta \vdash \sigma :_s \Gamma$}
\UnaryInfC{$\Gamma \vdash t[\sigma/] : A[\sigma/]$}
\DisplayProof
\quad (\textsc{back-inst})
\end{center}

\textsc{back-exact} applies when a prior vertex has a syntactically identical label. \textsc{back-inst} applies when a prior vertex matches under an explicit substitution $\sigma$. The judgement $\Delta \vdash \sigma :_s \Gamma$ states that $\sigma$ is an explicit substitution from context $\Gamma$ to $\Delta$, i.e., a list of terms in $\Delta$ that can replace the variables of $\Gamma$.

\subsection{Phases and focusing}

Rules are partitioned into two phases: an \textbf{asynchronous} phase (driving, deterministic) and a \textbf{synchronous} phase (split and backlink, requiring choice). This matches supercompilation: drive eagerly until stuck, then split or fold.

\subsection{Global progress condition}
\label{sec:global-condition}

A cyclic proof is sound if every infinite path contains infinitely many progress events (splits). A \emph{trace state} tracks the inductive variable being analyzed; a well-founded order requires strict decrease on at least one back-edge per cycle. This is enforced by the budget trace construction (Claim~2).

\subsection{Vertex classification}
\label{sec:vertex-classification}

Each vertex in the configuration graph is classified by its successor structure:

\begin{center}
\small
\begin{tabular}{c|c|c|c}
\textbf{Successors} & \textbf{Rule} & \textbf{SC operation} & \textbf{Phase} \\
\hline
$n = 0$ & \textsc{dr-leaf} & leaf (axiom) & — \\
$n = 1$, $\mathsf{drive\_cbn\_once}(t) \neq t$ & \textsc{dr-cbn-once} & driving (async) & invertible \\
$n = 1$, $\mathsf{drive\_cbn\_once}(t) = t$ & \textsc{dr-fold} & folding (memo hit) & cyclic closure \\
$n > 1$ & \textsc{dr-split} & case-split (sync) & choice point \\
\end{tabular}
\end{center}

\begin{proposition}[Vertex classification]
\label{prop:vertex-class}
Every vertex in the configuration graph produced by successful supercompilation satisfies exactly one of the four rules above.\footnote{Rocq: \texttt{cfg\_vertex\_drive\_rule} in \texttt{theories/Transform/SupercompilationCorrespondence.v}.}
\end{proposition}

\subsection{Cyclic proof structure (illustration)}
\label{sec:cyclic-diagram}

Figure~1 shows the cyclic proof for the length-map fusion example. The root goal (bottom) is $\mathsf{len}\,(\mathsf{map}\,f\,l) : \Nat$. Driving (rules \textsc{fix}, \textsc{$\iota$}, \textsc{comm}) reduces it to a case split on $l$. The nil branch trivially yields $0$. The cons branch reaches the recursive call $\mathsf{len}\,(\mathsf{map}\,f\,xs) : \Nat$, which is closed by a backlink to the root---the current goal is an instance of the root under the substitution $[l\mapsto xs]$, witnessed by the \textsc{back} rule. The split on $l$ ($\star$) is a progress event; the trace condition requires every cycle to contain one, ensuring structural descent.

\begin{center}
\footnotesize
\[
\AxiomC{}
\RightLabel{\;\fontsize{6}{7}\selectfont\textsc{nil}}
\UnaryInfC{$\Gamma \vdash 0 : \Nat$}
\qquad
\AxiomC{$\Gamma, xs \vdash [l\mapsto xs]:_s \Gamma$}
\RightLabel{\;\fontsize{6}{7}\selectfont\textsc{back} \textcolor{red!70}{\textdagger}}
\UnaryInfC{\textcolor{red!60}{$\Gamma, xs \vdash \mathsf{len}\,(\mathsf{map}\,f\,xs) : \Nat$ \textcolor{red!70}{\textdagger}}}
\RightLabel{\;\fontsize{6}{7}\selectfont\textsc{cons}}
\BinaryInfC{$\Gamma, l:\mathsf{List} \vdash \langle l \mid \mathit{nil}\to 0;\; \mathit{cons}\to\ldots \rangle : \Nat$}
\RightLabel{\;\fontsize{6}{7}\selectfont\textcolor{red!70}{\textsc{split} $\star$}}
\UnaryInfC{$\Gamma \vdash \langle l \mid \mathit{nil}\to 0;\; \mathit{cons}\to\ldots \rangle : \Nat$}
\RightLabel{\;\fontsize{6}{7}\selectfont\textsc{comm}}
\UnaryInfC{$\Gamma \vdash \mathsf{len}\,(\langle l \mid \ldots \rangle) : \Nat$}
\RightLabel{\;\fontsize{6}{7}\selectfont\textsc{$\iota$}}
\UnaryInfC{$\Gamma \vdash \mathsf{len}\,(\mathsf{map}\,f\,l) : \Nat$ \textcolor{red!70}{\textdagger}}
\RightLabel{\;\fontsize{6}{7}\selectfont\textsc{fix}}
\DisplayProof
\]

\vspace{-4pt}
{\fontsize{6}{7}\selectfont\textbf{Fig.\ 1.} Cyclic proof for length-map fusion. Red judgement = companion; $\dagger$ = backlink target; $\star$ = progress event.}
\end{center}

